% =====================================================================
\section{Phase 4: generating hypotheses with an LLM}\label{sec:p4}
% =====================================================================

\textbf{Goal.} Until now a human typed each candidate equation into the code. Phase~4 adds a generator that
proposes \emph{which} equations are worth testing, from a description of what was observed, without being told
the answer. The rule that makes this safe: \emph{LLM = hypothesis generator, PINN = validator.}

\textbf{The description the LLM sees.} \fn{build_problem_description} computes only statistics of the noisy
observations and turns them into text. For the Phase~1 data it says, among other things: the peak amplitude is
about 0.983 near $t=0.05$ and about 0.395 near $t=0.95$ (``decreased substantially''); the field stays near
zero at the edges; and the centroid moves only from $x\approx0.53$ to $0.41$ (``little evidence of bulk
translation''). It never names a mechanism.

\textbf{The key design choice: a generic equation compiler.} Phase~3 needed a hand-written Python function per
equation. That does not scale to whatever an LLM might propose. \fn{compile_residual} instead parses the
equation \emph{string} with SymPy (a symbolic-maths library) and compiles it into exactly the residual function
the Phase~3 trainer expects. For example, from \texttt{"u\_t = D*u\_xx + r*u*(1-u/K)"} it detects the fields
$u, u_t, u_{xx}$ and the parameters $D, r, K$ automatically. Any combination of $u, u_t, u_x, u_{xx}, u_{xxx}$
and named parameters works with zero code changes.

\textbf{Deterministic checks before any training} (\fn{validate_hypothesis}):
\begin{enumerate}
  \item \emph{Syntax}: can SymPy parse the equation?
  \item \emph{Variables}: is $u$ the dependent variable and are $x, t$ the independent ones?
  \item \emph{Parameters}: does every symbol used appear in the declared parameter list, and vice versa?
  \item \emph{Derivatives}: do the LLM's stated \fn{required_derivatives} match what SymPy finds?
  \item \emph{Units}: what fraction of parameters declare units (reported, not a proof)?
  \item \emph{Duplicates}: \fn{canonical_signature} replaces every parameter by a placeholder and ignores sign
        and side, so \texttt{u\_t = D*u\_xx} and \texttt{k*u\_xx = u\_t} are recognised as the same equation.
\end{enumerate}

\begin{figure}[!htb]
\centering
% flowchart 'p4' -- shared by project_guide.tex and submission_report.tex
\begin{tikzpicture}[node distance=3.6mm and 7mm]
  \node[fnw=cExp] (main) {06_hypothesis_generation.main};
  \node[fnw=cData, below=of main] (desc) {build_problem_description\\{\normalfont\tiny data statistics $\to$ plain text}};
  \node[fnw=cLlm, below=of desc] (gp) {get_provider\\{\normalfont\tiny config + env: gemini or mock}};
  \node[fnw=cHyp, below=of gp] (gen) {generate_hypotheses};
  \node[fnw=cLlm, right=of gen] (prompt) {build_hypothesis_prompt\\{\normalfont\tiny + SYSTEM\_PROMPT}};
  \node[fnw=cLlm, below=of prompt] (call) {provider.generate\\{\normalfont\tiny MockLLMProvider or GeminiProvider}};
  \node[fnw=cHyp, below=of gen, yshift=-12mm] (json) {_extract_json_array\\{\normalfont\tiny strip code fences, json.loads}};
  \node[fnw=cLlm, below=of json] (schema) {Hypothesis (pydantic)\\{\normalfont\tiny ParameterSpec; schema errors kept}};
  \node[fnw=cHyp, below=of schema] (val) {validate_hypothesis\\{\normalfont\tiny checks 1--5}};
  \node[fnw=cHyp, right=of val] (parse) {parse_equation\\{\normalfont\tiny SymPy: string $\to$ residual}};
  \node[fnw=cHyp, below=of val] (dedup) {deduplicate\\{\normalfont\tiny via canonical\_signature}};
  \node[fnw=cHyp, below=of dedup] (rank) {rank_hypotheses\\{\normalfont\tiny investigation order only}};
  \node[fnw=cHyp, right=of rank] (scores) {_completeness_score\\_simplicity_score};
  \node[fnw=cHyp, below=of rank] (adapt) {hypothesis_to_pinn_config};
  \node[fnw=cHyp, right=of adapt] (comp) {compile_residual\\{\normalfont\tiny SymPy lambdify $\to$ torch}};
  \node[fnw=cPinn, below=of adapt] (train) {build_and_train\\{\normalfont\tiny the unchanged Phase 3 trainer}};
  \node[fnw=cVal, below=of train] (score) {build_validation_result};
  \draw[flow] (main) -- (desc); \draw[flow] (desc) -- (gp); \draw[flow] (gp) -- (gen);
  \draw[flow] (gen) -- (prompt); \draw[flow] (prompt) -- (call);
  \draw[flow] (call) |- node[note, above, pos=0.75]{raw text} (json);
  \draw[flow] (json) -- (schema); \draw[flow] (schema) -- (val); \draw[flow] (val) -- (parse);
  \draw[flow] (val) -- (dedup); \draw[flow] (dedup) -- (rank); \draw[flow] (rank) -- (scores);
  \draw[dflow] (rank) -- node[note, right]{optional PINN run} (adapt);
  \draw[flow] (adapt) -- (comp); \draw[flow] (adapt) -- (train); \draw[flow] (train) -- (score);
\end{tikzpicture}

\pkglegend
\caption{Phase~4 function flow: describe $\to$ prompt $\to$ LLM $\to$ parse JSON $\to$ schema $\to$ deterministic
checks $\to$ deduplicate $\to$ rank $\to$ (optionally) compile into a PINN and validate.}
\label{fig:p4}
\end{figure}

\textbf{Results (mock LLM, all four candidates run through the Phase~3 PINN).}
\begin{center}
\small
\begin{tabular}{@{}llll@{}}
\toprule
\textbf{ID} & \textbf{Equation} & \textbf{PINN status} & \textbf{Fitted parameters} \\ \midrule
H1 & $u_t = D\,u_{xx}$ & supported & $D=0.09989$ \\
H2 & $u_t + c\,u_x = 0$ & rejected & $c=0.0164$ \\
H3 & $u_t = D\,u_{xx} + r\,u(1-u/K)$ & supported & $D=0.1006,\ r=0.0148,\ K=1.215$ \\
H4 & $u_t + c\,u_x = D\,u_{xx}$ & supported & $c=-0.0004,\ D=0.0999$ \\
\bottomrule
\end{tabular}
\end{center}

\begin{lesson}[: validation alone cannot reject harmless extra terms]
H3 and H4 \emph{contain} diffusion as a special case. Their extra parameters converged to almost exactly the
``switch this term off'' values ($r\approx0.015$, $c\approx0$), so they fit as well as H1 and pass the same
thresholds. A pass/fail validator cannot tell ``minimal correct model'' from ``needlessly general model that
also fits''. This is the \textbf{nested-model problem} (Section~\ref{sec:selection}), and it is why Phase~6 adds
term necessity, complexity and BIC to the selection.
\end{lesson}

% =====================================================================
\section{Phase 5: scientific RAG and a knowledge graph}\label{sec:p5}
% =====================================================================

\textbf{Goal.} Ground hypothesis generation in retrieved scientific evidence (RAG), and organise that evidence as
a knowledge graph with provenance, so every fact can be traced to its source.

\textbf{The corpus.} Five short Markdown documents in \fn{rag/corpus/}: diffusion, advection, reaction,
reaction--diffusion and advection--diffusion. Each has a front-matter header (id, title, source, domain,
declared mechanism, declared equations, declared parameters) and textbook-style sections.

\textbf{Preventing answer leakage.} Documents state only the generic textbook law (e.g.\ ``Fick's second law,
$u_t=D\,u_{xx}$''). They never state this dataset's $D=0.1$ and never say ``this data is diffusion''. A real
scientist knows the textbook candidates before looking at new data; knowing the answer would be cheating. An
automated check confirms the string ``0.1'' appears in no corpus document.

\begin{figure}[!htb]
\centering
% flowchart 'p5' -- shared by project_guide.tex and submission_report.tex
\begin{tikzpicture}[node distance=3.6mm and 6mm]
  \node[fnw=cExp] (main) {07_scientific_rag_kg.main};
  \node[fnw=cRag, below=of main] (lc) {load_corpus\\{\normalfont\tiny read the 5 files + front matter}};
  % RAG column (left)
  \node[fnw=cRag, below left=6mm and -10mm of lc] (ing) {RAGPipeline.ingest};
  \node[fnw=cRag, below=of ing] (chunk) {chunk_document\\{\normalfont\tiny split at section headers}};
  \node[fnw=cRag, below=of chunk] (emb) {LocalEmbeddingProvider.embed\\{\normalfont\tiny word hashing $\to$ 256-dim unit vector}};
  \node[fnw=cRag, below=of emb] (vs) {InMemoryVectorStore.add};
  % KG column (right)
  \node[fnw=cKg, below right=6mm and -10mm of lc] (bg) {build_graph_from_documents};
  \node[fnw=cKg, below=of bg] (ent) {KnowledgeGraph.add_entity\\{\normalfont\tiny Document, Mechanism, Variable, \ldots}};
  \node[fnw=cKg, below=of ent] (rel) {KnowledgeGraph.add_relationship\\{\normalfont\tiny governs, has-parameter, \ldots}};
  \node[fnw=cKg, below=of rel] (src) {_find_chunk_id\\{\normalfont\tiny provenance: which chunk said it}};
  % merge
  \node[fnw=cData, below=12mm of vs] (pd) {build_problem_description\\{\normalfont\tiny the query text}};
  \node[fnw=cRag, below=of pd, xshift=27mm] (ctx) {build_scientific_context};
  \node[fnw=cRag, below left=5mm and -12mm of ctx] (ret) {Retriever.retrieve\\{\normalfont\tiny cosine top-k, metadata filters}};
  \node[fnw=cRag, below=of ret] (ev) {evidence_from_retrieval};
  \node[fnw=cKg, below right=5mm and -12mm of ctx] (q) {equations_for_mechanism\\parameters_for_mechanism\\variables_for_mechanism};
  \node[fnw=cRag, below=31mm of ctx] (sc) {ScientificContext\\{\normalfont\tiny mechanisms, equations, parameters, sources}};
  \node[fnw=cHyp, below=of sc] (gen) {generate_hypotheses(context)\\{\normalfont\tiny Mode B: evidence-grounded prompt}};
  \node[fnw=cKg, right=of lc, xshift=4mm] (ser) {save_graph\\{\normalfont\tiny graph to JSON (results/rag\_kg)}};
  \draw[flow] (main) -- (lc); \draw[flow] (lc) -- (ing); \draw[flow] (lc) -- (bg);
  \draw[flow] (ing) -- (chunk); \draw[flow] (chunk) -- (emb); \draw[flow] (emb) -- (vs);
  \draw[flow] (bg) -- (ent); \draw[flow] (ent) -- (rel); \draw[flow] (rel) -- (src);
  \draw[dflow] (chunk.east) -- (bg.west |- chunk.east);
  \draw[dflow] (bg.north) |- (ser.west);
  \draw[flow] (vs) -- (pd); \draw[flow] (pd) -- (ctx); \draw[flow] (src) |- (ctx);
  \draw[flow] (ctx) -- (ret); \draw[flow] (ret) -- (ev); \draw[flow] (ctx) -- (q);
  \draw[flow] (ev) |- (sc); \draw[flow] (q) |- (sc); \draw[flow] (sc) -- (gen);
\end{tikzpicture}

\pkglegend
\caption{Phase~5 function flow. Left: RAG indexing (chunk, embed, store). Right: knowledge-graph construction
with provenance (the dashed arrow: chunks give each relationship its source). Bottom: the observation
description is used as a query; retrieved evidence and graph queries merge into one \fn{ScientificContext} that
is passed to the hypothesis generator.}
\label{fig:p5}
\end{figure}

\textbf{Results.}
\begin{itemize}
  \item The graph has \textbf{22 entities} (5 documents, 5 mechanisms, 3 variables, 4 parameters, 5 equations)
        and \textbf{40 relationships, all sourced}. Shared concepts merge: there is a single node for $D$,
        even though three documents declare it, and repeated facts merge their sources instead of creating
        duplicate edges.
  \item Example queries: asking \fn{equations_for_mechanism} about diffusion returns \texttt{u\_t = D*u\_xx},
        and asking \fn{parameters_for_mechanism} returns \texttt{D}.
  \item For the Phase~1 data, the context contained the mechanisms advection, diffusion and reaction, their
        three textbook equations, and the parameters $D, K, c, r$.
  \item \textbf{Mode A} (no context) and \textbf{Mode B} (with context) produced the identical four candidate equations. That is
        expected: the mock LLM returns canned answers whatever the prompt. Showing that evidence changes the
        proposals needs a real LLM (Section~\ref{sec:gemini}).
\end{itemize}

\begin{lesson}[: RAG and the graph do different jobs]
RAG answers ``which text is relevant?'' by similarity; it is fuzzy and ranked. The graph answers ``how are these
documented concepts related, and who said so?''; it is exact and traceable. Neither decides what is true.
Phase~7 measures their real contribution and finds it modest: RAG gives 100\% recall of the right mechanism on a
5-document corpus, but top-1 retrieval is right only 1 time in 6.
\end{lesson}
